\section{Exact Computer-Assisted Certificate for the Mixed Cap Overlaps}
\label{sec:strategy4-exact-certificate}

The two mixed cap overlaps in Strategy~4 are certified by exact symbolic
arithmetic.  This section states the complete mathematical reduction, explains
the verification algorithms, and formally incorporates the code and sparse
integer data listed below as an electronic proof supplement.  No floating-point
arithmetic, interval arithmetic, adaptive subdivision, branch-and-bound, or
unrecorded numerical tolerance is used.

\subsection{Certificate manifest}

All paths in this subsection are relative to
\begin{verbatim}
proof/3XXX_CE0/31XX_Nplus1/310X_all_Vd0/
3105X_self_contained_direct_Vd0_nine_point/3105X_computation/
\end{verbatim}
The exact files incorporated in the certificate are
\begin{center}
\scriptsize
\renewcommand{\arraystretch}{1.08}
\begin{tabular}{p{.53\textwidth}p{.34\textwidth}}
\toprule
file&Git blob SHA\\
\midrule
\texttt{mixed\_overlap\_core\_data\_00.py}&
\texttt{359b219f9deda8586217dcbf9ccc10de0ad48578}\\
\texttt{mixed\_overlap\_core\_data\_01.py}&
\texttt{bfc807f906359636035634f799d74d552cf847fc}\\
\texttt{mixed\_overlap\_core\_data\_02.py}&
\texttt{92887696ef94aeb93f3d95ab9b16a5388406fff4}\\
\texttt{mixed\_overlap\_core\_data\_03.py}&
\texttt{c2791a8a6add132ad7a17655cbb13987124cb24e}\\
\texttt{mixed\_overlap\_core\_data\_04.py}&
\texttt{7eeabecbbd17543b728130a1f2bc84f79e11a8d2}\\
\texttt{mixed\_overlap\_core\_data\_05.py}&
\texttt{d676114292e722b60d1b792f615ee88ad58b6982}\\
\texttt{mixed\_overlap\_core\_polynomials.py}&
\texttt{07f57b178cf74980f8d956541842d0c4c4b2faa7}\\
\texttt{verify\_mixed\_overlap\_core\_derivation.py}&
\texttt{0a55211ae37f56e1fc461560b0668a01456c7417}\\
\texttt{verify\_global\_core\_positivity.py}&
\texttt{821f9a6dc750fbdb80e02e40282b32cc845c3383}\\
\bottomrule
\end{tabular}
\end{center}
The six data shards decode to one canonical sparse-polynomial transcript with
SHA--256 digest
\begin{verbatim}
dc46aaf263655d5159ecd3a81db72ee82477951d06172f4743b248df37209485
\end{verbatim}
The first verifier independently derives the transcript from the witness
formulas and compares every coefficient with this authenticated object.  The
second verifier derives all Bernstein coefficients from that same object and
checks their signs exactly.

\subsection{Rational radial upper envelopes}

Retain the notation of Section~\ref{sec:ab-core}.  Thus
\[
 p=1-b,\qquad q=1-a,\qquad s=p+q,\qquad
 m=\min\{p,q\},\qquad M=\max\{p,q\},
\]
\[
 \chi=s^4-s^2+mM,\qquad h=\frac{\sqrt3}{2},\qquad
 c_*=c_{\max}(p,q).
\]
The strict witness domain gives
\begin{equation}
 mM>(1-s)(2-s),\qquad m>1-s,\qquad s>\frac23.
 \label{eq:certificate-domain}
\end{equation}
Put
\[
 F_m(x)=x^4-x^2+mx-m^2,\qquad
 \kappa=F_m'(s)=4s^3-2s+m.
\]
Then
\[
 \kappa>4s^3-3s+1=(s+1)(2s-1)^2>0.
\]
Define
\begin{equation}
 \bar c_L=s-\frac{\chi}{\kappa}\quad(\chi\le0),
 \qquad
 \bar c_T=s-\frac{\chi}{1-m}\quad(\chi\ge0).
 \label{eq:certificate-cbar}
\end{equation}

\begin{lemma}[Branchwise radial upper bounds]
\label{lem:paper-branchwise-cbar}
On the corresponding radial cells,
\[
 c_*\le\bar c_L<1,\qquad c_*\le\bar c_T<1.
\]
At $\chi=0$ both upper bounds equal $s=c_*$.
\end{lemma}

\begin{proof}
On the $L$ cell, $F_m(s)=\chi\le0$ and $F_m(c_*)=0$.  For $x\ge s>2/3$,
$F_m''(x)>0$, so convexity gives
\[
 0\ge F_m(s)+F_m'(s)(c_*-s),
\]
and hence $c_*\le\bar c_L$.  If $U=1-s$, then
\[
 \chi+U\kappa
 >U\{m+1-4U+8U^2-3U^3\}>0
\]
by \eqref{eq:certificate-domain}; therefore $\bar c_L<1$.

On the $T$ cell put
\[
 E=\sqrt{4s^2-3},\qquad
 m_0=\frac{s(1-E)}2,\qquad M_0=\frac{s(1+E)}2.
\]
Then
\[
 \chi=(m-m_0)(M_0-m),\qquad
 s-c_*=\frac{2(m-m_0)}{1+E}.
\]
Consequently
\[
 \frac{\chi}{s-c_*}
 =\frac{1+E}{2}(M_0-m)\le1-m,
\]
which gives $c_*\le\bar c_T$.  The strict inequality $\chi>0$ gives
$\bar c_T<s<1$.
\end{proof}

Reflect so that $z=a-b\ge0$ and put $U=a+b-1$.  Then
\[
 s=1-U,\qquad m=\frac{1-U-z}{2},\qquad
 D^2=z^2+6U+3U^2=:K(U,z),
\]
and
\[
 \Xi(U,z)=4\chi
 =1-10U+21U^2-16U^3+4U^4-z^2,
\]
\[
 G(U,z)=2\kappa=5-21U+24U^2-8U^3-z.
\]
Thus
\begin{equation}
 \bar c_L=1-U-\frac{\Xi(U,z)}{2G(U,z)},
 \qquad
 \bar c_T=1-U-\frac{\Xi(U,z)}{2(1+U+z)}.
 \label{eq:certificate-rational-cbar}
\end{equation}

\subsection{Gram reduction and the geometric implication}

On the active cell let $\bar c$ denote the corresponding expression in
\eqref{eq:certificate-rational-cbar}, and put
\[
 \bar\eta=h(1-\bar c),\qquad e=1-\bar c,\qquad t=2\bar c-1.
\]
Lemma~\ref{lem:paper-branchwise-cbar} gives
$0<\bar\eta\le\eta=h(1-c_*)$.  It therefore suffices to prove the mixed
intersections for the smaller disks with radii $2\bar\eta,\bar\eta,2\bar\eta$.

Put $C'=\mathsf R^{-1}C$ and write
\[
 N_A=\lVert A\rVert^2,\qquad N_C=\lVert C\rVert^2,\qquad
 \nu=\langle A,C'\rangle,\qquad
 \Delta=\operatorname{cross}(C',A)>0.
\]
The Gram identity is
\[
 \Delta^2=N_AN_C-\nu^2.
\]
Define
\begin{equation}
 Q_A=\Delta^2-\lVert tA-eC'\rVert^2,\qquad
 Q_C=\Delta^2-\lVert tC'-eA\rVert^2.
 \label{eq:certificate-Q}
\end{equation}

\begin{lemma}[Residual signs imply the mixed cap intersections]
\label{lem:paper-residual-to-overlap}
If $Q_A,Q_C\ge0$, then
\[
 I(C,2\bar\eta)\cap I(C+\mathsf RA,\bar\eta)\ne\varnothing,
\]
and
\[
 I(C+\mathsf RA,\bar\eta)\cap I(\mathsf RA,2\bar\eta)\ne\varnothing.
\]
Consequently the same intersections hold with $\eta$ in place of
$\bar\eta$.
\end{lemma}

\begin{proof}
Rotate the first pair by $\mathsf R^{-1}$.  Its centers become $C'$ and
$A+C'$, so their difference is $A$.  An outer common tangent normal $n_A$ is
chosen with
\[
 \langle A,n_A\rangle=\bar\eta.
\]
Since $\lVert A\rVert>\bar\eta$ and $\Delta>0$, choose the sign for which
\[
 \langle C',n_A\rangle
 =\frac{\bar\eta\nu+\sqrt{N_A-\bar\eta^2}\,\Delta}{N_A}.
\]
The common support of the two disks is then
\[
 2\bar\eta+
 \frac{\bar\eta\nu+\sqrt{N_A-\bar\eta^2}\,\Delta}{N_A}.
\]
Let $\bar\tau=h-2\bar\eta=ht$.  Expanding
\eqref{eq:certificate-Q} and using the Gram identity gives
\begin{equation}
 h^2N_AQ_A
 =(N_A-\bar\eta^2)\Delta^2
 -(\bar\tau N_A-\bar\eta\nu)^2.
 \label{eq:certificate-PA}
\end{equation}
Thus $Q_A\ge0$ implies
\[
 \sqrt{N_A-\bar\eta^2}\,\Delta
 \ge\bar\tau N_A-\bar\eta\nu,
\]
and the displayed common support is at least
$2\bar\eta+\bar\tau=h$.  Hence the two level-$h$ caps intersect.
The same calculation with $A$ and $C'$ exchanged gives
\begin{equation}
 h^2N_CQ_C
 =(N_C-\bar\eta^2)\Delta^2
 -(\bar\tau N_C-\bar\eta\nu)^2,
 \label{eq:certificate-PC}
\end{equation}
and proves the second intersection.  Enlarging the disk radii from
$\bar\eta$ to $\eta$ only enlarges the caps.
\end{proof}

\subsection{Eight exact integer-polynomial signs}

The Newton witnesses $A,C$ are rational functions of $U,z,D$.  Substituting
\eqref{eq:certificate-rational-cbar} into \eqref{eq:certificate-Q}, clearing
denominators, and reducing by
\[
 D^2-K(U,z)=0
\]
gives, for $B\in\{L,T\}$ and $X\in\{A,C\}$,
\begin{equation}
 \operatorname{num}(Q_{B,X})
 =32(K+3)^2
 \{A_{B,X}(U,z)+D B_{B,X}(U,z)\},
 \label{eq:certificate-core-reduction}
\end{equation}
where all eight core polynomials lie in $\mathbb Z[U,z]$.
Every omitted denominator is strictly positive on the strict domain: its
factors are positive constants, $(D+3)^4$, $G(U,z)^2$ or
$(1+U+z)^2$, and squares of the two Newton-derivative numerators.  Here
$G=2\kappa>0$, and the fixed-line sign theorem proves both Newton derivatives
strictly negative.

The derivation verifier constructs $A,C$, both rational envelopes, and all four
residuals directly in the exact field $\mathbb Q(U,z,D)$.  It performs the
reduction modulo $D^2-K$, removes the common positive factor, primitive-normalizes
the eight integer polynomials, and compares their sparse term lists
coefficient by coefficient with the authenticated transcript.  Thus the data
are not an independently trusted precomputation; they are checked against the
geometric formulas.

\subsection{Three global positive-basis certificates}

Set
\[
 \mathcal G(U)=1-6U-3U^2,\qquad
 H(U)=1-10U+21U^2-16U^3+4U^4,
\]
\[
 U_H=1-\frac{\sqrt3}{2},\qquad
 U_{\max}=\frac2{\sqrt3}-1.
\]
The complete reflected domain is the union of the three closed cells
\[
\begin{array}{ll}
T:&0\le U\le U_H,\quad 0\le z^2\le H(U),\\
L_0:&0\le U\le U_H,\quad H(U)\le z^2\le\mathcal G(U),\\
L_1:&U_H\le U\le U_{\max},\quad0\le z^2\le\mathcal G(U).
\end{array}
\]
For the $T$ cell use
\[
 U=\frac{(x-1)(3-x)}{4x},\qquad
 z=y\frac{9-x^4}{8x^2},\qquad
 1\le x\le\sqrt3,\quad0\le y\le1,
\]
and then $x=1+(\sqrt3-1)r$.  This maps $[0,1]^2$ onto the complete cell.
After multiplication by positive powers of $2$ and $x$, the four transformed
core polynomials have global tensor Bernstein expansions with data
\[
\begin{array}{c|ccc}
&\text{bidegree}&\text{coefficient count}&\text{zero count}\\ \hline
A_{T,A}&(88,22)&2047&2\\
B_{T,A}&(80,20)&1701&0\\
A_{T,C}&(88,22)&2047&2\\
B_{T,C}&(80,20)&1701&0.
\end{array}
\]
Every nonzero coefficient is strictly positive.

For an $L$-cell polynomial write
\[
 F(U,z)=E_F(U,w)+zO_F(U,w),\qquad w=z^2,
\]
and
\[
 R_F(U,w)=E_F(U,w)^2-wO_F(U,w)^2.
\]
The implications $E_F\ge0$ and $R_F\ge0$ give
$E_F\ge\sqrt w\,|O_F|$, hence $F\ge0$ for $z\ge0$.
Use the two exact charts
\[
 U=U_Hr,\qquad
 w=H(U)+s\{\mathcal G(U)-H(U)\}
\]
for $L_0$, and
\[
 U=U_H+(U_{\max}-U_H)r,\qquad
 w=s\mathcal G(U)
\]
for $L_1$.  The positivity verifier forms one global Bernstein expansion for
each of $E_F,R_F$ on each square.  The resulting sixteen expansions have the
bidegrees and coefficient counts recorded in
Appendix~\ref{subsec:exact-mixed-overlap}; every coefficient is strictly
positive.

The verifier implements polynomial addition, multiplication, substitution,
and power-to-Bernstein conversion over exact integers.  Coefficients in
$\mathbb Q(\sqrt3)$ are represented as rational pairs and signed by exact
rational comparisons.  In particular, a polynomial
$\sum a_{ij}r^is^j$ of bidegree $(m,n)$ is converted using
\[
 b_{k\ell}
 =\sum_{i\le k}\sum_{j\le\ell}
 a_{ij}\frac{\binom{k}{i}}{\binom{m}{i}}
       \frac{\binom{\ell}{j}}{\binom{n}{j}},
\]
which is the exact tensor Bernstein coefficient formula.

\begin{theorem}[Exact mixed-overlap certificate]
\label{thm:paper-exact-mixed-certificate}
On the complete strict witness domain with $c_*>2/3$, all eight polynomials in
\eqref{eq:certificate-core-reduction} are nonnegative.  Hence
$Q_A,Q_C\ge0$, and the two mixed cap intersections required by
Proposition~\ref{prop:technical-four-overlaps} hold.
\end{theorem}

\begin{proof}
The $T$ chart proves the four $T$-cell core signs by nonnegative Bernstein
coefficients.  On each $L$ chart the positive Bernstein certificates for
$E_F$ and $R_F$ prove the four $L$-cell signs.  Since $D\ge0$,
\eqref{eq:certificate-core-reduction} gives $Q_A,Q_C\ge0$ for $z\ge0$;
reflection supplies $z\le0$.  Lemma~\ref{lem:paper-residual-to-overlap} then
gives the two cap intersections.
\end{proof}

\begin{remark}[Reproduction status]
The files above constitute the exact certificate used by the proof.  A
reproduction run should execute
\begin{verbatim}
verify_mixed_overlap_core_derivation.py
verify_global_core_positivity.py
\end{verbatim}
from their package directory.  This manuscript revision does not claim that
those commands were rerun during the present source-only edit; the certificate
is incorporated by its exact repository objects and authenticated transcript.
\end{remark}
